\section{基础构件：增强型 LLM}

在介绍具体的 Workflow 模式之前，需要先理解 agentic 系统的基础构件——\textbf{Augmented LLM}（增强型 LLM）。

增强型 LLM 是在基础 LLM 之上叠加了三种核心能力的模块：

\begin{table}[H]
\centering
\begin{tabular}{>{\bfseries}l p{10cm}}
\toprule
增强能力 & 说明 \\
\midrule
Retrieval（检索） & 从外部知识库、数据库或文档中获取相关信息，扩展模型的知识边界 \\
Tools（工具） & 调用外部 API、执行代码、操作文件等，使模型能够影响外部世界 \\
Memory（记忆） & 在多轮交互中保持上下文，记录和回忆之前的信息 \\
\bottomrule
\end{tabular}
\caption{增强型 LLM 的三大核心能力}
\end{table}

当前的模型已能\textbf{主动}使用这些能力——自行生成搜索查询、选择合适的工具、决定保留哪些信息。Anthropic 建议在实现时关注两个方面：

\begin{enumerate}
    \item \textbf{针对具体用例定制}这些能力，而非追求通用性。
    \item 确保这些能力为 LLM 提供\textbf{简洁、文档完备的接口}。
\end{enumerate}

\begin{knowledgebox}{Model Context Protocol (MCP)}
Anthropic 推出的 MCP 协议允许开发者通过统一的客户端实现接入不断增长的第三方工具生态。这意味着增强型 LLM 的工具层可以标准化、可复用，而非每个项目重新开发。MCP 正成为 Agent 工具生态的重要基础设施。
\end{knowledgebox}

\subsection{本章小结}
增强型 LLM 是所有 agentic 系统的基础模块，由 Retrieval、Tools、Memory 三种能力增强。后续所有 Workflow 和 Agent 模式都假设每次 LLM 调用都具备这些增强能力。


